\documentclass[11pt,a4paper]{article}
\usepackage[utf8]{inputenc}
\usepackage[margin=1in]{geometry}
\usepackage{amsmath,amssymb}
\usepackage{booktabs}
\usepackage{xcolor}
\usepackage{hyperref}
\usepackage{fancyhdr}
\usepackage{enumitem}

\setlength{\headheight}{14pt}

\hypersetup{
    colorlinks=true,
    linkcolor=blue,
    filecolor=magenta,      
    urlcolor=cyan,
    pdftitle={Project Report: Brushless DC (BLDC) Motor Driver Circuit Using Sensorless ESC Architecture},
    pdfpagemode=UseOutlines,
}

\pagestyle{fancy}
\fancyhf{}
\rhead{Brushless DC}
\lhead{Technical Project Report}
\rfoot{Page \thepage}

\title{\textbf{\Huge Engineering Project Report}\\[0.5em] \Large Project 093: Brushless DC (BLDC) Motor Driver Circuit Using Sensorless ESC Architecture}
\author{\textbf{Bigyanlabs Engineering Team}\\ Electronics \& Embedded Systems Division}
\date{\today}

\begin{document}

\maketitle
\thispagestyle{empty}

\begin{abstract}
This technical report presents the comprehensive design, circuit schematics, bill of materials, mathematical modeling, assembly methodology, and experimental testing for \textbf{Brushless DC (BLDC) Motor Driver Circuit Using Sensorless ESC Architecture}. Classified under the \textbf{Advanced / Expert (Engineering Level)} tier in the \textbf{Drone / Avionics ESC} category, this design provides optimized performance, high reliability, and clear practical utility for school and engineering applications.
\end{abstract}

\newpage
\tableofcontents
\newpage

\section{Introduction \& System Overview}
The objective of this project is to implement a robust electronic system for \textbf{Brushless DC (BLDC) Motor Driver Circuit Using Sensorless ESC Architecture}. This document details the step-by-step engineering design, component functions, mathematical principles, and experimental results.

\section{Bill of Materials (Components List)}
\begin{table}[h!]
\centering
\caption{Bill of Materials for Brushless DC (BLDC) Motor Driver Circuit Using Sensorless ESC Architecture}
\vspace{0.5em}
\begin{tabular}{clcp{7cm}}
\toprule
\textbf{S.No.} & \textbf{Component Name} & \textbf{Qty} & \textbf{Circuit Function} \\
\midrule
1 & Power Transistors / High-End MCU & 2 & Power MOSFETs / IGBTs / STM32 / ESP32 \\
2 & High-Frequency Transformers & 1 & Ferro-magnetic energy storage core \\
3 & Gate Drivers / Fast Diodes & 2 & IR2110 gate driver and Ultrafast diodes \\
4 & Sensing Amplifiers / CT Sensors & 2 & Current sensors (ACS712/SCT-013) and Differential amps \\
5 & Filter Capacitors (Low ESR) & 2 & Low ESR electrolytic and Polypropylene film caps \\
6 & Auxiliary Power \\& Protection & 1 Set & Auxiliary SMPS, heatsinks, and snubber networks \\
\bottomrule
\end{tabular}
\end{table}

\section{System Architecture \& Methodology}
The Brushless DC (BLDC) Motor Driver Circuit Using Sensorless ESC Architecture is an advanced engineering implementation featuring high-frequency switching, closed-loop feedback algorithms, or high-speed DSP processing for maximum efficiency and precision.

\section{Circuit Assembly \& Testing Procedure}
\begin{enumerate}
    \item Place core active components and microcontrollers/ICs on the breadboard or PCB.
    \item Wire DC supply rails (+VCC and GND) with appropriate decoupling capacitors.
    \item Connect sensor networks and input signal conditioning circuits.
    \item Interface output switches, MOSFET drivers, or visual/acoustic indicators.
    \item Apply power and measure operating node voltages to verify proper performance.
\end{enumerate}

\section{Mathematical Derivations \& Governing Equations}
\begin{equation}
V_{\text{out}} = \frac{V_{\text{in}}}{1 - D}
\end{equation}
\begin{equation}
\eta = \frac{P_{\text{out}}}{P_{\text{in}}} \times 100\%
\end{equation}
\begin{equation}
f_r = \frac{1}{2\pi \sqrt{L C}}
\end{equation}

\section{Applications \& Engineering Conclusion}
Industrial automation, renewable energy systems, electric vehicle power trains, aerospace electronics, and high-performance drone / avionics esc.

\end{document}
